\section{Implementação}

Os registros selecionados documentam desenvolvimento e operação em pomares, com participantes e atividades identificados, mas não disponibilizam dados experimentais com amostragem e comparação controlada. A Tabela~\ref{tab:evidencias_implementacao} delimita as conclusões sustentadas.

\begin{table*}[!htb]
    \centering
    \caption{Evidências de implementação e alcance dos resultados.}
    \label{tab:evidencias_implementacao}
    \small
    \setlength{\tabcolsep}{5pt}
    \renewcommand{\arraystretch}{1.10}

    \begin{tabularx}{\textwidth}{@{}
        >{\raggedright\arraybackslash}p{0.24\textwidth}
        >{\raggedright\arraybackslash}p{0.42\textwidth}
        >{\raggedright\arraybackslash}X
        @{}}
        \hline
        \textbf{Registro} &
        \textbf{Evidência disponível} &
        \textbf{Conclusão sustentada} \\
        \hline

        Rivoira, Itália, 2021~\cite{tevel2021rivoira} &
        Conclusão relatada do primeiro piloto comercial nos pomares do grupo. &
        Ensaio em ambiente real; resultados qualitativos declarados. \\
        \hline

        Unifrutti, Chile, 2023~\cite{unifrutti2023tevel} &
        Colheita de maçãs com oito robôs e envio ao beneficiamento. &
        Implantação com período e configuração identificados. \\
        \hline

    \end{tabularx}
\end{table*}

O comunicado sobre a Rivoira relata poucas contusões, classificação satisfatória por cor e alta cobertura de colheita~\cite{tevel2021rivoira}. Segundo Maor, o piloto operou continuamente por cinco semanas, com comparação diária entre frutos colhidos por robôs e trabalhadores. As contusões, inicialmente semelhantes às da colheita manual, tornaram-se inferiores após uma correção na unidade terrestre. Também foi declarada diferença favorável superior a 10\% na classificação, sem definição da métrica. O ensaio teve escala limitada, sem divulgação do total de maçãs colhidas, da proporção acessível aos robôs ou de indicadores consolidados de rentabilidade~\cite{ellis2022tevel}.

A Unifrutti informa implantação entre março e maio de 2023 e análise dos frutos no beneficiamento, com resultados de seletividade e minimização de contusões. Não apresenta produtividade, taxa de danos, grupo comparador ou horas de intervenção humana. O intervalo informado não comprova operação ininterrupta~\cite{unifrutti2023tevel}.

Os relatos sustentam a operação em pomares e registram aperfeiçoamentos técnicos~\cite{tevel2021rivoira,ellis2022tevel,unifrutti2023tevel}. A comparação com a colheita manual limita-se a danos e classificação seletiva~\cite{ellis2022tevel}. Faltam dados para determinar ganhos de produtividade, relacionar o número de robôs à vazão em toneladas por hora ou quantificar a redução de custos operacionais.

O prêmio divulgado pela Kubota em 2022 reconheceu um conceito desenvolvido com a Tevel, combinando \acp{FAR}, tratores autônomos e outras máquinas agrícolas~\cite{kubota2022agrifuture}. O reconhecimento contextualiza o desenvolvimento e a integração tecnológica e industrial, sem quantificar desempenho. A evidência de implementação deste estudo provém das operações em campo; premiações e anúncios de parceria documentam propostas e iniciativas de integração.

A operação diurna e noturna anunciada pela Tevel~\cite{tevel2026} expressa uma capacidade declarada. A disponibilidade efetiva exige contabilizar interrupções por manutenção, reposicionamento, manejo dos recipientes, falhas e condições ambientais. Sem esses registros, ``24/7'' não constitui uma medição de disponibilidade.